\section{The Elements of ALBERT}
\label{sec:model}
In this section, we present the design decisions for ALBERT and provide quantified comparisons against corresponding configurations of the original BERT architecture~\citep{devlin2018bert}.
% Studies are now in the next section.
% as well as careful ablation studies that shed light on the impact of these design decisions.


\subsection{Model architecture choices}
The backbone of the ALBERT architecture is similar to BERT in that it uses a transformer encoder~\citep{vaswani2017attention} with GELU nonlinearities~\citep{hendrycks2016gelus}. 
We follow the BERT notation conventions and denote the vocabulary embedding size as $E$, the number of encoder layers as $L$, and the hidden size as $H$. Following \cite{devlin2018bert}, 
we set the feed-forward/filter size to be $4H$ and the number of attention heads to be $H/64$. 

There are three main contributions that ALBERT makes over the design choices of BERT.

\paragraph{Factorized embedding parameterization.}
In BERT, as well as subsequent modeling improvements such as \xlnet~\citep{yang2019xlnet} and \roberta~\citep{liu2019roberta}, the WordPiece embedding size $E$ is tied with the hidden layer size $H$, i.e., $E\equiv H$.
This decision appears suboptimal for both modeling and practical reasons, as follows.

From a modeling perspective, WordPiece embeddings are meant to learn {\em context-independent} representations, whereas hidden-layer embeddings are meant to learn {\em context-dependent} representations.
As experiments with context length indicate~\citep{liu2019roberta}, the power of BERT-like representations comes from the use of context to provide the signal for learning such context-dependent representations.
As such, untying the WordPiece embedding size $E$ from the hidden layer size $H$ allows us to make a more efficient usage of the total model parameters as informed by modeling needs, which dictate that $H \gg E$. 

From a practical perspective, natural language processing usually require the vocabulary size $V$ to be large.\footnote{Similar to BERT, all the experiments in this paper use a vocabulary size $V$ of 30,000.} 
If $E\equiv H$, then increasing $H$ increases the size of the embedding matrix, which has size $V \times E$.
This can easily result in a model with billions of parameters, most of which are only updated sparsely during training. 

Therefore, for ALBERT we use a factorization of the embedding parameters, decomposing them into two smaller matrices.
Instead of projecting the one-hot vectors directly into the hidden space of size $H$, we first project them into a lower dimensional embedding space of size $E$, and then project it to the hidden space.
By using this decomposition, we reduce the embedding parameters from $O(V \times H)$  to $O(V \times E + E \times H)$.
This parameter reduction is significant when $H \gg E$. We choose to use the same E for all word pieces because they are much more evenly distributed across documents compared to whole-word embedding, where having different embedding size  (\cite{grave2017efficient, baevski2018adaptive,dai2019transformer} ) for different words is important.  


\paragraph{Cross-layer parameter sharing.}
For ALBERT, we propose cross-layer parameter sharing as another way to improve parameter efficiency.
There are multiple ways to share parameters, e.g., only sharing feed-forward network (FFN) parameters across layers, or only sharing attention parameters.
The default decision for ALBERT is to share all parameters across layers. All our experiments use this default decision unless otherwise specified.
We compare this design decision against other strategies in our experiments in Sec.~\ref{sec:sharing}.

Similar strategies have been explored by \cite{dehghani2018universal} (Universal Transformer, UT) and \cite{bai2019deep} (Deep Equilibrium Models, DQE) for Transformer networks.
Different from our observations, \cite{dehghani2018universal} show that UT outperforms a vanilla Transformer.
% even without dynamical halting, % RS: not a must here, and introduces a new concept that's not relevant
\cite{bai2019deep} show that their DQEs reach an equilibrium point for which the input and output embedding of a certain layer stay the same.
Our measurement on the L2 distances and cosine similarity show that our embeddings are oscillating rather than converging.


\begin{figure}[!htbp] 
\begin{subfigure}{.5\textwidth}
  \centering
  \includegraphics[width=.8\linewidth]{adjacent_l2}
  \label{fig:adjacent_l2}
\end{subfigure}%
\begin{subfigure}{.5\textwidth}
  \centering
  \includegraphics[width=.8\linewidth]{adjacent_cos_sim}
  \label{fig:adjacent_cos_sim}
\end{subfigure}
\caption{The L2 distances and cosine similarity (in terms of degree) of the input and output embedding of each layer for BERT-large and ALBERT-large.}
\label{fig:l2_cos_sim}
\end{figure}

Figure \ref{fig:l2_cos_sim} shows the L2 distances and cosine similarity of the input and output embeddings for each layer, using BERT-large and ALBERT-large configurations (see Table~\ref{table:model_architecture}).
We observe that the transitions from layer to layer are much smoother for ALBERT than for BERT.
These results show that weight-sharing has an effect on stabilizing network parameters.
Although there is a drop for both metrics compared to BERT, they nevertheless do not converge to 0 even after 24 layers.
This shows that the solution space for ALBERT parameters is very different from the one found by DQE. 

\paragraph{Inter-sentence coherence loss.}
In addition to the masked language modeling (MLM) loss~\citep{devlin2018bert}, BERT uses an additional loss called next-sentence prediction (NSP).
NSP is a binary classification loss for predicting whether two segments appear consecutively in the original text, as follows:
positive examples are created by taking consecutive segments from the training corpus;
negative examples are created by pairing segments from different documents;
positive and negative examples are sampled with equal probability.
The NSP objective was designed to improve performance on downstream tasks, such as natural language inference, that require reasoning about the relationship between sentence pairs. 
However, subsequent studies~\citep{yang2019xlnet,liu2019roberta} found NSP's impact unreliable and decided to eliminate it, a decision supported by an improvement in downstream task performance across several tasks.

We conjecture that the main reason behind NSP's ineffectiveness is its lack of difficulty as a task, as compared to MLM.
As formulated, NSP conflates {\em topic prediction} and {\em coherence prediction} in a single task\footnote{Since a negative example is constructed using material from a different document, the negative-example segment is misaligned both from a topic and from a coherence perspective.}.
However, topic prediction is easier to learn compared to coherence prediction, and also overlaps more with what is learned using the MLM loss.
%RS: the second reason not supported for now by the empirical results in Table~\ref{table:sop}.
% Second, for MLM, negative examples only have half of the context of the positive examples, which makes it difficult for the model to learn to use larger context. 

%For ALBERT, we maintain that inter-sentence modeling is an important aspect of language understanding, if defined as a {\em coherence} loss. 
We maintain that inter-sentence modeling is an important aspect of language understanding, but we propose a loss based primarily on {\em coherence}. 
That is, for ALBERT, we use a sentence-order prediction (SOP) loss, which avoids topic prediction and instead focuses on modeling inter-sentence coherence.
The SOP loss uses as positive examples the same technique as BERT (two consecutive segments from the same document), and as negative examples the same two consecutive segments but with their order swapped.
This forces the model to learn finer-grained distinctions about discourse-level coherence properties.
As we show in Sec.~\ref{sec:sop}, it turns out that NSP cannot solve the SOP task at all (i.e., it ends up learning the easier topic-prediction signal, and performs at random-baseline level on the SOP task), while SOP can solve the NSP task to a reasonable degree, presumably based on analyzing misaligned coherence cues. 
As a result, ALBERT models consistently improve downstream task performance for multi-sentence encoding tasks.

\subsection{Model setup}
We present the differences between BERT and ALBERT models with comparable hyperparameter settings in Table~\ref{table:model_architecture}.
Due to the design choices discussed above, ALBERT models have much smaller parameter size compared to corresponding BERT models.


\begin{table}
\small
\centering
\begin{tabular}{l  l | c c c c c c}
 \multicolumn{2}{c}{Model}          &Parameters	&Layers &Hidden &Embedding  &Parameter-sharing \\
\hline
\multirow{3}{*}{BERT}    & base	    &108M    &12 & 768  & 768  & False\\
                         & large	&334M	 &24 & 1024 & 1024 & False\\
 %                        & xlarge	&1270M	 &24 & 2048 & 2048 & False\\
    \hline
\multirow{4}{*}{ALBERT}  &base	    &12M     &12 & 768  & 128  & True \\
                         &large	    &18M	 &24 & 1024 & 128  & True\\
                         &xlarge	&60M	 &24 & 2048 & 128  & True\\
                         &xxlarge	&235M	 &12 & 4096 & 128  & True\\
\end{tabular}
\caption{The configurations of the main BERT and ALBERT models analyzed in this paper.}
\label{table:model_architecture}
\end{table}

For example, ALBERT-large has about 18x fewer parameters compared to BERT-large, 18M versus 334M.
% If we set BERT to have an extra-large size with $H=2048$, we end up with a model that has 1.27 billion parameters and the model becomes difficult to run without model parallization.  and under-performs (Fig.~\ref{fig:large_vs_extra_large}).
An ALBERT-xlarge configuration with $H=2048$ has only 60M parameters and an ALBERT-xxlarge configuration with $H=4096$ has 233M parameters, i.e., around 70\% of BERT-large's parameters.
Note that for ALBERT-xxlarge, we mainly report results on a 12-layer network because a 24-layer network (with the same configuration) obtains similar results but is computationally more expensive.

This improvement in parameter efficiency is the most important advantage of ALBERT's design choices. 
Before we can quantify this advantage, we need to introduce our experimental setup in more detail.


\section{Experimental Results}

\subsection{Experimental Setup}
\label{sec:setup}
To keep the comparison as meaningful as possible, we follow the \bert~\citep{devlin2018bert} setup in using the \textsc{BookCorpus}~\citep{zhu2015aligning} and English Wikipedia~\citep{devlin2018bert} for pretraining baseline models.
These two corpora consist of around 16GB of uncompressed text.
We format our inputs as ``\textsc{[CLS]} $x_1$ \textsc{[SEP]} $x_2$ \textsc{[SEP]}'', where $x_1=x_{1,1},x_{1,2}\cdots$ and $x_2=x_{1,1},x_{1,2}\cdots$ are two segments.\footnote{A segment is usually comprised of more than one natural sentence, which has been shown to benefit performance by~\cite{liu2019roberta}.} 
We always limit the maximum input length to 512, and randomly generate input sequences shorter than 512 with a probability of 10\%.
Like \bert, we use a vocabulary size of 30,000, tokenized using SentencePiece~\citep{kudo-richardson-2018-sentencepiece} as in \xlnet~\citep{yang2019xlnet}.

We generate masked inputs for the MLM targets using $n$-gram masking~\citep{joshi2019spanbert}, with the length of each $n$-gram mask selected randomly.
The probability for the length $n$ is given by
$$
    p(n) = \frac{1/n}{\sum_{k=1}^N{1/k}}
$$
We set the maximum length of $n$-gram (i.e., $n$) to be 3 (i.e., the MLM target can consist of up to a 3-gram of complete words, such as ``White House correspondents'').

All the model updates use a batch size of 4096 and a \textsc{Lamb} optimizer with learning rate 0.00176~\citep{you2019reducing}.
We train all models for 125,000 steps unless otherwise specified.
Training was done on Cloud TPU V3.
The number of TPUs used for training ranged from 64 to 512, depending on model size. 

The experimental setup described in this section is used for all of our own versions of BERT as well as ALBERT models, unless otherwise specified. 


\subsection{Evaluation Benchmarks}


\subsubsection{Intrinsic Evaluation}
% RS: this needs to be more detailed.
To monitor the training progress, we create a development set based on the development sets from \squad and RACE using the same procedure as in Sec.~\ref{sec:setup}. 
We report accuracies for both MLM and sentence classification tasks. Note that we only use this set to check how the model is converging; it has not been used in a way that would affect the performance of any downstream evaluation, such as via model selection. 
%\lanzhzh{add a clarification to make sure that people would not think that we misuse the downstream datasets.}\kevin{thanks for the clarification!}


\subsubsection{Downstream Evaluation}
Following \cite{yang2019xlnet} and \cite{liu2019roberta}, we evaluate our models on three popular benchmarks: The General Language Understanding Evaluation (GLUE) benchmark~\citep{wang-etal-2018-glue}, two versions of the Stanford Question Answering Dataset (\squad;~\citealp{rajpurkar-etal-2016-squad,rajpurkar-etal-2018-know}), and the ReAding Comprehension from Examinations  (RACE) dataset~\citep{lai2017race}. 
For completeness, we provide description of these benchmarks in Appendix~\ref{downstream_detailed_description}.
As in \citep{liu2019roberta}, we perform early stopping on the development sets, on which we report all comparisons except for our final comparisons based on the task leaderboards, for which we also report test set results. For GLUE datasets that have large variances on the dev set, we report median over 5 runs.  


\subsection{Overall Comparison between BERT and ALBERT}
We are now ready to quantify the impact of the design choices described in Sec.~\ref{sec:model}, specifically the ones around parameter efficiency.
The improvement in parameter efficiency showcases the most important advantage of ALBERT's design choices, as shown in Table~\ref{table:overall_comparison}: with only around 70\% of BERT-large's parameters, ALBERT-xxlarge achieves significant improvements over BERT-large, as measured by the difference on development set scores for several representative downstream tasks:
SQuAD v1.1 (+1.9\%), SQuAD v2.0 (+3.1\%), MNLI (+1.4\%), SST-2 (+2.2\%), and RACE (+8.4\%).


% We also observe that BERT-xlarge gets significantly worse results than BERT-base on all metrics.
% This indicates that a model like BERT-xlarge is more difficult to train than those that have smaller parameter sizes. 
Another interesting observation is the speed of data throughput at training time under the same training configuration (same number of TPUs). 
Because of less communication and fewer computations, ALBERT models have higher data throughput compared to their corresponding BERT models.
If we use BERT-large as the baseline, we observe that ALBERT-large is about 1.7 times faster in iterating through the data while ALBERT-xxlarge is about 3 times slower because of the larger structure. 
%The 12-layer ALBERT-xxlarge model outperform all other models significantly.
%Compared to BERT-large, which has roughly 1.45x more parameters, ALBERT-xxlarge has an average of $3.6\%$ absolute improvements over 5 downstream tasks and $8.5\%$ absolute improvements over the RACE task. 


\begin{table}[!htbp] 
\setlength{\tabcolsep}{5pt}
\small
\centering
\begin{tabular}{ l l | c c c c c c | c | c}
 \multicolumn{2}{c}{Model}    &Parameters        	&SQuAD1.1	&SQuAD2.0	&MNLI	&SST-2	&RACE	&Avg & Speedup\\
\hline
\multirow{3}{*}{BERT} &base	     &108M       &90.4/83.2	&80.4/77.6	&84.5	&92.8	&68.2	&82.3 & 4.7x \\

    &large	        &334M	&92.2/85.5	&85.0/82.2	&86.6	&93.0	&73.9	&85.2 & 1.0\\
%    &xlarge			&1270M &86.4/78.1	&75.5/72.6	&81.6	&90.7	&54.3	&76.6 & 1.0\\
\hline 
\multirow{4}{*}{ALBERT}  &base	      &12M    	    &89.3/82.3	&80.0/77.1	&81.6	&90.3	&64.0	&80.1 &5.6x \\
&large	        	 &18M   &90.6/83.9	&82.3/79.4	&83.5	&91.7	&68.5	&82.4 &1.7x\\
&xlarge		   &60M   &92.5/86.1	&86.1/83.1	&86.4	&92.4	&74.8	&85.5 & 0.6x\\
&xxlarge	&235M	& \textbf{94.1/88.3}	&\textbf{88.1/85.1}	&\textbf{88.0}	&\textbf{95.2}	&\textbf{82.3}	&\textbf{88.7} & 0.3x\\
\end{tabular}
\caption{Dev set results for models pretrained over \textsc{BookCorpus} and Wikipedia for 125k steps.
Here and everywhere else, the Avg column is computed by averaging the scores of the downstream tasks to its left (the two numbers of F1 and EM for each SQuAD are first averaged).}
\label{table:overall_comparison}
\end{table}

Next, we perform ablation experiments that quantify the individual contribution of each of the design choices for ALBERT. 

\subsection{Factorized Embedding Parameterization}

%\subsubsection{Effect of vocabulary embedding size}
Table \ref{table:voc_embedding} shows the effect of changing the vocabulary embedding size $E$ using an ALBERT-base configuration setting (see Table~\ref{table:model_architecture}), using the same set of representative downstream tasks.
Under the non-shared condition (BERT-style), larger embedding sizes give better performance, but not by much.
Under the all-shared condition (ALBERT-style), an embedding of size 128 appears to be the best.
Based on these results, we use an embedding size $E=128$ in all future settings, as a necessary step to do further scaling.

\begin{table}[!htbp] 
\small
\centering
\begin{tabular}{ c c c | c c c c c | c }
Model  & $E$ & Parameters     	&SQuAD1.1	&SQuAD2.0	&MNLI	&SST-2	&RACE	&Avg \\
\hline
\multirow{4}{*}{\shortstack{ALBERT \\ base \\ not-shared}}
& 64  & 87M &89.9/82.9	&80.1/77.8	&82.9	&91.5	&66.7   &81.3 \\
& 128 & 89M &89.9/82.8	&80.3/77.3	&83.7	&91.5	&67.9	&81.7 \\
& 256 & 93M &90.2/83.2	&80.3/77.4	&84.1	&91.9	&67.3	&81.8 \\
& 768 & 108M &90.4/83.2	&80.4/77.6	&84.5	&92.8	&68.2	&82.3  \\
\hline
\multirow{4}{*}{\shortstack{ALBERT \\base \\ all-shared}}
&64     & 10M &88.7/81.4	&77.5/74.8  &80.8	&89.4	&63.5	&79.0 \\
&128    & 12M  &89.3/82.3	&80.0/77.1	&81.6	&90.3	&64.0	&80.1 \\
&256    & 16M &88.8/81.5	&79.1/76.3	&81.5	&90.3	&63.4	&79.6 \\
&768    & 31M &88.6/81.5	&79.2/76.6	&82.0	&90.6	&63.3	&79.8 \\
\end{tabular}
\caption{The effect of vocabulary embedding size on the performance of ALBERT-base.}
\label{table:voc_embedding}
\end{table}

% RS: I propose we take this out, since we're not following this trend.
% As a side observation with respect to model capacity, for ALBERT-xxlarge, we can reduce the number of parameters for the word-piece embedding from 61M, which is more than half of the overall parameters, to 4M. 

\vspace{-3pt}

\subsection{Cross-layer parameter sharing}
\label{sec:sharing}
Table \ref{table:paramter-sharing} presents experiments for various cross-layer parameter-sharing strategies, using an ALBERT-base configuration (Table~\ref{table:model_architecture}) with two embedding sizes ($E=768$ and $E=128$).
We compare the all-shared strategy (ALBERT-style), the not-shared strategy (BERT-style), and intermediate strategies in which only the attention parameters are shared (but not the FNN ones) or only the FFN parameters are shared (but not the attention ones).

The all-shared strategy hurts performance under both conditions, but it is less severe for $E=128$ (-1.5 on Avg) compared to $E=768$ (-2.5 on Avg).
In addition, most of the performance drop appears to come from sharing the FFN-layer parameters, while sharing the attention parameters results in no drop when $E=128$ (+0.1 on Avg), and a slight drop when $E=768$ (-0.7 on Avg).

There are other strategies of sharing the parameters cross layers. For example, We can divide the $L$ layers into $N$ groups of size $M$, and each size-$M$ group shares parameters. Overall, our experimental results shows that the smaller the group size $M$ is, the better the performance we get. However, decreasing group size $M$ also dramatically increase the number of overall parameters. We choose all-shared strategy as our default choice.  

\begin{table}[!htbp] 
\small
\centering
\begin{tabular}{ l l | c | c c c c c | c }
\multicolumn{2}{c}{Model}	    &Parameters  	&SQuAD1.1	&SQuAD2.0	&MNLI	&SST-2	&RACE &Avg \\
\hline
\multirow{4}{*}{\shortstack{ALBERT \\ base \\ $E$=768}}
&all-shared       &31M   &88.6/81.5	&79.2/76.6   &82.0	&90.6	&63.3	&79.8 \\
&shared-attention &83M   &89.9/82.7	&80.0/77.2	 &84.0	&91.4	&67.7	&81.6 \\
&shared-FFN       &57M   &89.2/82.1	&78.2/75.4	 &81.5	&90.8	&62.6	&79.5 \\
&not-shared        &108M  &90.4/83.2	&80.4/77.6	 &84.5	&92.8	&68.2	&82.3 \\
\hline
\multirow{4}{*}{\shortstack{ALBERT \\base \\ $E$=128}}
&all-shared       & 12M  &89.3/82.3	&80.0/77.1	&82.0	&90.3	&64.0	&80.1 \\
&shared-attention & 64M  &89.9/82.8	&80.7/77.9	&83.4	&91.9	&67.6	&81.7 \\
&shared-FFN       & 38M  &88.9/81.6	&78.6/75.6	&82.3	&91.7	&64.4	&80.2 \\
&not-shared   & 89M  &89.9/82.8	&80.3/77.3	&83.2	&91.5	&67.9	&81.6 \\
\end{tabular}
\caption{The effect of cross-layer parameter-sharing strategies, ALBERT-base configuration.}
\label{table:paramter-sharing}
\end{table}

\vspace{-5pt}

\subsection{Sentence order prediction (SOP)}
\label{sec:sop}
We compare head-to-head three experimental conditions for the additional inter-sentence loss: none (XLNet- and RoBERTa-style), NSP (BERT-style), and SOP (ALBERT-style), using an ALBERT-base configuration.
Results are shown in Table \ref{table:sop}, both over intrinsic (accuracy for the MLM, NSP, and SOP tasks) and downstream tasks.

\begin{table}[!htbp] 
\small
\centering
\begin{tabular}{ l | c c c | c c c c  c | c }
&\multicolumn{3}{|c|}{Intrinsic Tasks} & \multicolumn{6}{c}{Downstream Tasks} \\
SP tasks	    	&MLM	&NSP & SOP &SQuAD1.1	&SQuAD2.0	&MNLI	&SST-2	&RACE & Avg \\
\hline
None &54.9 &52.4 & 53.3 &88.6/81.5	&78.1/75.3	&81.5	&89.9	&61.7		& 79.0 \\
NSP	 &54.5 &90.5 & 52.0 &88.4/81.5	&77.2/74.6	&81.6	&\textbf{91.1}	&62.3	& 79.2  \\
SOP	 &54.0 &78.9 &86.5  &\textbf{89.3/82.3}	&\textbf{80.0/77.1}	&\textbf{82.0}	&90.3	&\textbf{64.0}	 & \textbf{80.1}	\\
\end{tabular}
\caption{The effect of sentence-prediction loss, NSP vs. SOP, on intrinsic and downstream tasks.}
\label{table:sop}
\end{table}

The results on the intrinsic tasks reveal that the NSP loss brings no discriminative power to the SOP task (52.0\% accuracy, similar to the random-guess performance for the ``None'' condition).
This allows us to conclude that NSP ends up modeling only topic shift.
In contrast, the SOP loss does solve the NSP task relatively well (78.9\% accuracy), and the SOP task even better (86.5\% accuracy).
Even more importantly, the SOP loss appears to consistently improve downstream task performance for multi-sentence encoding tasks (around +1\% for SQuAD1.1, +2\% for SQuAD2.0, +1.7\% for RACE), for an Avg score improvement of around +1\%.

\subsection{What if we train for the same amount of time?}

The speed-up results in Table \ref{table:overall_comparison} indicate that data-throughput for BERT-large is about 3.17x higher compared to ALBERT-xxlarge.
Since longer training usually leads to better performance, we perform a comparison in which, instead of controlling for data throughput (number of training steps), we control for the actual training time (i.e., let the models train for the same number of hours).
In Table \ref{table:same-training-time}, we compare the performance of a  BERT-large model after 400k training steps (after 34h of training), roughly equivalent with the amount of time needed to train an ALBERT-xxlarge model with 125k training steps (32h of training).


\begin{table}[!htbp] 
\small
\centering
\begin{tabular}{ l c c|  c c c c c | c }
Models	& Steps& Time	&SQuAD1.1	&SQuAD2.0	&MNLI	&SST-2	&RACE & Avg \\
\hline
BERT-large & 400k  & 34h      &93.5/87.4	&86.9/84.3	&87.8	&94.6	&77.3	&87.2\\
ALBERT-xxlarge & 125k & 32h  &\textbf{94.0/88.1}	&\textbf{88.3/85.3}	&87.8	&\textbf{95.4}	&\textbf{82.5}   &\textbf{88.7}\\
\end{tabular}
\caption{The effect of controlling for training time, BERT-large vs ALBERT-xxlarge configurations.}
\label{table:same-training-time}
\end{table}

After training for roughly the same amount of time, ALBERT-xxlarge is significantly better than BERT-large: +1.5\% better on Avg, with the difference on RACE as high as +5.2\%.


\subsection{Additional training data and dropout effects}
The experiments done up to this point use only the Wikipedia and \textsc{BookCorpus} datasets, as in \citep{devlin2018bert}.
In this section, we report measurements on the impact of the additional data used by both XLNet~\citep{yang2019xlnet} and RoBERTa~\citep{liu2019roberta}.

Fig.~\ref{fig:additiona_data} plots the dev set MLM accuracy under two conditions, without and with additional data, with the latter condition giving a significant boost.
We also observe  performance improvements on the downstream tasks in Table \ref{table:additional_training_data}, except for the SQuAD benchmarks (which are Wikipedia-based, and therefore are negatively affected by out-of-domain training material).

\begin{figure}[!htbp] 
\begin{subfigure}{.5\textwidth}
  \centering
  \includegraphics[width=.8\linewidth]{adding_additional_data}
  \caption{Adding data   \label{fig:additiona_data}}
\end{subfigure}%
\begin{subfigure}{.5\textwidth}
  \centering
  \includegraphics[width=.8\linewidth]{removing_dropout}
  \caption{Removing dropout   \label{fig:removing_dropout}}
\end{subfigure}
\caption{The effects of adding data and removing dropout during training.}
\label{fig:additional_data_removing_dropout}
\end{figure}

\begin{table}[!htbp] 
\small
\centering
\begin{tabular}{ r |  c c c c c | c }
%Number of layer		
&SQuAD1.1	&SQuAD2.0	&MNLI	&SST-2	&RACE & Avg \\
\hline
No additional data      &\textbf{89.3/82.3}	&\textbf{80.0/77.1}	&81.6	&90.3	&64.0	&80.1\\
With additional data      &88.8/81.7	&79.1/76.3	&\textbf{82.4}	&\textbf{92.8}	&\textbf{66.0}	&\textbf{80.8}\\
\end{tabular}
\caption{The effect of additional training data using the ALBERT-base configuration.}
\label{table:additional_training_data}
\end{table}


We also note that, even after training for 1M steps, our largest models still do not overfit to their training data.
As a result, we decide to remove dropout to further increase our model capacity.
The plot in Fig.~\ref{fig:removing_dropout} shows that removing dropout significantly improves MLM accuracy.
Intermediate evaluation on ALBERT-xxlarge at around 1M training steps (Table~\ref{table:removing_dropout}) also confirms that removing dropout helps the downstream tasks.
There is empirical \citep{szegedy2017inception} and theoretical \citep{li2019understanding} evidence showing that a combination of batch normalization and dropout in Convolutional Neural Networks may have harmful results.
To the best of our knowledge, we are the first to show that dropout can hurt performance in large Transformer-based models. However, the underlying network structure of ALBERT is a special case of the transformer and further experimentation is needed to see if this phenomenon appears with other transformer-based architectures or not.

\begin{table}[!htbp] 
\small
\centering
\begin{tabular}{ r |  c c c c c | c }
%Number of layer		
&SQuAD1.1	&SQuAD2.0	&MNLI	&SST-2	&RACE & Avg \\
\hline
With dropout     &94.7/89.2	&89.6/86.9	&90.0	&96.3	&85.7	&90.4\\
Without dropout  &\textbf{94.8/89.5}	 &\textbf{89.9/87.2}	&\textbf{90.4}	&\textbf{96.5}	&\textbf{86.1}	&\textbf{90.7}\\
\end{tabular}
\caption{The effect of removing dropout, measured for an ALBERT-xxlarge configuration.}
\label{table:removing_dropout}
\end{table}


\vspace{-5pt}
%\radu{I moved two subsections with ablation experiments to the Appendix.}
%\lanzhzh{cool, look good}

\subsection{Current State-of-the-art on NLU Tasks}

The results we report in this section make use of the training data used by~\cite{devlin2018bert}, as well as the additional data used by~\cite{liu2019roberta} and~\cite{yang2019xlnet}. 
We report state-of-the-art results under two settings for fine-tuning: single-model and ensembles.
In both settings, we only do single-task fine-tuning\footnote{Following \cite{liu2019roberta}, we fine-tune for RTE, STS, and MRPC using an MNLI checkpoint.}.
Following \cite{liu2019roberta}, on the development set we report the median result over five runs.

%\mingda{Added this} 


\begin{table}[!htbp] 
\small
\centering
\begin{tabular}{ l   c c c c c c c c c c}
Models &MNLI	&QNLI	&QQP	&RTE	&SST	&MRPC	&CoLA	&STS	&WNLI	&Avg \\
\hline
\multicolumn{11}{l}{\textit{Single-task single models on dev}}\\
BERT-large	 &86.6	&92.3	&91.3	&70.4	&93.2	&88.0	&60.6	&90.0 &- &-	\\	
XLNet-large	 &89.8	&93.9	&91.8	&83.8	&95.6	&89.2	&63.6	&91.8 &- &-	\\	
RoBERTa-large&90.2	&94.7	&\textbf{92.2}	&86.6	&96.4	&\textbf{90.9}	&68.0	&92.4 &- &-	\\	
ALBERT (1M)   &90.4	&95.2	&92.0	&88.1	&96.8	&90.2	&68.7	&92.7 &- &-	\\
ALBERT (1.5M) &\textbf{90.8}	&\textbf{95.3}	&\textbf{92.2}	&\textbf{89.2}	&\textbf{96.9}	&\textbf{90.9}	&\textbf{71.4}	&\textbf{93.0} &- &-	\\	
\hline
\multicolumn{11}{l}{\textit{Ensembles on test (from leaderboard as of Sept. 16, 2019)}}\\
ALICE	    &88.2	&95.7	&\textbf{90.7}	&83.5	&95.2	&92.6	&\textbf{69.2}	&91.1	&80.8	&87.0 \\
MT-DNN	    &87.9	&96.0	&89.9	&86.3	&96.5	&92.7	&68.4	&91.1	&89.0	&87.6 \\
XLNet	    &90.2	&98.6	&90.3	&86.3	&96.8	&93.0	&67.8	&91.6	&90.4	&88.4 \\
RoBERTa	    &90.8	&98.9	&90.2	&88.2	&96.7	&92.3	&67.8	&92.2	&89.0	&88.5 \\
Adv-RoBERTa	&91.1	&98.8	&90.3	&88.7	&96.8	&93.1	&68.0	&92.4	&89.0	&88.8 \\
ALBERT 	    &\textbf{91.3}	&\textbf{99.2}	&90.5	&\textbf{89.2}	&\textbf{97.1}	&\textbf{93.4}	&69.1	&\textbf{92.5}	&\textbf{91.8}	&\textbf{89.4} \\
\end{tabular}
\caption{State-of-the-art results on the GLUE benchmark.
For single-task single-model results, we report ALBERT at 1M steps (comparable to RoBERTa) and at 1.5M steps. The ALBERT ensemble uses models trained with 1M, 1.5M, and other numbers of steps.}
\label{table:glue}
\end{table}


The single-model ALBERT configuration incorporates the best-performing settings discussed: an ALBERT-xxlarge configuration (Table~\ref{table:model_architecture}) using combined MLM and SOP losses, and no dropout. 
The checkpoints that contribute to the final ensemble model are selected based on development set performance; the number of checkpoints considered for this selection range from 6 to 17, depending on the task. 
For the GLUE (Table~\ref{table:glue}) and RACE (Table~\ref{table:sota-squad-race}) benchmarks, we average the model predictions for the ensemble models, where the candidates are fine-tuned from different training steps using the 12-layer and 24-layer architectures.
For SQuAD (Table~\ref{table:sota-squad-race}), we average the prediction scores for those spans that have multiple probabilities; we also average the scores of the ``unanswerable'' decision.

Both single-model and ensemble results indicate that ALBERT improves the state-of-the-art significantly for all three benchmarks, achieving a GLUE score of 89.4, a SQuAD 2.0 test F1 score of 92.2, and a RACE test accuracy of 89.4. 
The latter appears to be a particularly strong improvement, a jump of +17.4\% absolute points over BERT~\citep{devlin2018bert,clark2019bam}, +7.6\% over XLNet~\citep{yang2019xlnet},  +6.2\% over RoBERTa~\citep{liu2019roberta}, and 5.3\% over DCMI+ \citep{zhang2019dcmn+}, an ensemble of multiple models specifically designed for reading comprehension tasks.
Our single model achieves an accuracy of $86.5\%$, which is still $2.4\%$ better than the state-of-the-art ensemble model. 



\begin{table}[!htbp] 
\setlength{\tabcolsep}{4pt}
\small
\centering
\begin{tabular}{ l   c c c c }
Models		&SQuAD1.1 dev	&SQuAD2.0 dev	& SQuAD2.0 test &RACE test (Middle/High) \\
\hline
\multicolumn{5}{l}{\textit{Single model (from leaderboard as of Sept.~23, 2019)}}\\
BERT-large          & 90.9/84.1 & 81.8/79.0 & 89.1/86.3 & 72.0 (76.6/70.1) \\
XLNet               & 94.5/89.0 & 88.8/86.1 & 89.1/86.3 & 81.8 (85.5/80.2) \\
RoBERTa             & 94.6/88.9 &89.4/86.5  & 89.8/86.8 & 83.2 (86.5/81.3) \\
UPM                 &  - & - & 89.9/87.2 & - \\
XLNet + SG-Net Verifier++ & - & - & 90.1/87.2 &  -\\
ALBERT (1M)          & 94.8/89.2 & 89.9/87.2 & - & 86.0 (88.2/85.1) \\
ALBERT (1.5M)        & \textbf{94.8/89.3}	& \textbf{90.2/87.4} & \textbf{90.9/88.1} & \textbf{86.5 (89.0/85.5)} \\
\hline
\multicolumn{5}{l}{\textit{Ensembles (from leaderboard as of Sept.~23, 2019)}}\\
BERT-large          & 92.2/86.2        & - & - &  - \\
XLNet + SG-Net Verifier & -& -& 	90.7/88.2	& -\\
UPM & - & - & 90.7/88.2 & \\ 
XLNet + DAAF + Verifier & - & - & 	90.9/88.6 & -\\
DCMN+               & - & -& - & 84.1 (88.5/82.3) \\
ALBERT     &    \textbf{95.5/90.1}      & \textbf{91.4/88.9} & \textbf{92.2/89.7}    &\textbf{89.4 (91.2/88.6)}  \\ 
\end{tabular}
\caption{State-of-the-art results on the SQuAD and RACE benchmarks. }
\label{table:sota-squad-race}
\vspace{-3pt}
\end{table}
	

